\section{Introduction}
\label{sec:introduction}

In urban mixed-traffic autonomous driving, planners built on path--velocity decomposition emit unnecessary fallback stops because time-blind path bounds and obstacle-wise side choices discard interaction timing before either solver runs. For autonomous-vehicle system integrators and safety engineers, these artificial deadlocks directly degrade ride quality and route efficiency, creating a measurable gap between theoretical feasibility and operational performance in dense traffic. This paper develops MIKU, a constraint reconstruction framework for two-stage quadratic programming path--velocity decomposition architectures, evaluated on 4,000 deterministic interval instances and 3,500 paired simulated scenarios via collision-free goal reaching, progress ratio, jerk RMS, and planning-time quantiles, with degraded-stop rate directly reflecting the operational trade-offs among safety, progress, ride comfort, and real-time computation \citep{gonzalez2016,paden2016,schwarting2018planning}.

Specifically, path--velocity decomposition separates lateral path generation in Frenet station--lateral coordinates $s$--$l$ from longitudinal speed optimization in station--time coordinates $s$--$t$ \citep{kant1986pvd,werling2010,zhou2021dliaps}. Industrial implementations typically solve two structured convex quadratic programs that expose environmental constraints through spatial path boundaries and temporal ST boundaries \citep{fan2018baidu,zhang2020pathqp,wang2022_jfr}. In dense dynamic scenes this separation discards two kinds of information: a path boundary from a current snapshot or a horizon union cannot distinguish an area occupied now from one free at the ego's arrival, and irreversible obstacle-wise side choices can raise the lower and lower the upper path bound in mutually incompatible ways even when a jointly selected passage remains open.

While joint space--time search, mixed-integer optimization, and spatio-temporal corridor methods coordinate spatial and temporal decisions simultaneously \citep{mcnaughton2011lattice,kessler2023,yoon2024stcorridor}, they typically expand state spaces or alter optimization structures, making them incompatible with production two-QP architectures. Consequently, current industrial PVD implementations suffer from three fundamental limitations:
first, time-blind path boundaries merge occupancies across the prediction horizon, losing arrival-conditioned clearance before lateral optimization;
second, greedy, obstacle-wise side choices cannot capture group-level continuous passages, causing lower and upper path bounds to cross and falsely signal a blockage even when viable corridors exist;
and third, existing cross-stage coordination mechanisms cannot transfer spatio-temporal homotopy decisions without modifying downstream convex solvers.
As illustrated in Figure~\ref{fig:teaser_en}, MIKU resolves these challenges directly at the constraint interface by constructing time-dependent lateral bands, selecting spatial and temporal homotopy classes, and mapping them into standard box constraints accepted by the two QPs.

\begin{figure}
\centering
\adjustbox{max width=\textwidth}{\begin{tikzpicture}[x=0.82cm,y=0.82cm,>=Stealth]
\begin{scope}
  \fill[bglight] (-.4,-2.25) rectangle (7.8,2.25);
  \draw[gray,thick] (-.4,2.25)--(7.8,2.25) (-.4,-2.25)--(7.8,-2.25);
  \draw[dashed,gray!60] (-.4,0)--(7.8,0);
  \fill[dangerred!70] (2.7,.45) rectangle (4.3,1.0);
  \fill[dangerred!70] (2.7,-1.0) rectangle (4.3,-.45);
  \draw[dangerred!60,dashed] (2.3,.05) rectangle (4.7,1.4);
  \draw[dangerred!60,dashed] (2.3,-1.4) rectangle (4.7,-.05);
  \node[font=\scriptsize] at (3.5,.73) {$O_a$};
  \node[font=\scriptsize] at (3.5,-.73) {$O_b$};
  \node[deeppink,font=\scriptsize\bfseries] at (5.0,.8) {$R$};
  \node[deepblue,font=\scriptsize\bfseries] at (2.0,-.8) {$L$};
  \draw[deeppink,very thick] (-.4,2.25)--(2.1,2.25)--(2.1,0)--(4.9,0)--(4.9,2.25)--(7.8,2.25);
  \draw[deepblue,very thick] (-.4,-2.25)--(2.1,-2.25)--(2.1,0)--(4.9,0)--(4.9,-2.25)--(7.8,-2.25);
  \node[dangerred,font=\scriptsize\bfseries] at (3.5,1.75) {$l^+\le l^-\;\Rightarrow\;$ blocked};
  \node[font=\small\bfseries] at (3.5,-2.75) {Obstacle-wise greedy choice};
\end{scope}
\begin{scope}[xshift=9cm]
  \fill[bglight] (-.4,-2.25) rectangle (7.8,2.25);
  \draw[gray,thick] (-.4,2.25)--(7.8,2.25) (-.4,-2.25)--(7.8,-2.25);
  \draw[dashed,gray!60] (-.4,0)--(7.8,0);
  \fill[dangerred!70] (2.7,.45) rectangle (4.3,1.0);
  \fill[dangerred!70] (2.7,-1.0) rectangle (4.3,-.45);
  \draw[dangerred!60,dashed] (2.3,.05) rectangle (4.7,1.4);
  \draw[dangerred!60,dashed] (2.3,-1.4) rectangle (4.7,-.05);
  \node[font=\scriptsize] at (3.5,.73) {$O_a$};
  \node[font=\scriptsize] at (3.5,-.73) {$O_b$};
  \node[deepblue,font=\scriptsize\bfseries] at (5.0,.8) {$R$};
  \node[deepblue,font=\scriptsize\bfseries] at (5.0,-.8) {$R$};
  \fill[lightgreen,opacity=.28] (-.4,-2.25) rectangle (7.8,-1.45);
  \draw[deeppink,very thick] (-.4,2.25)--(2.1,2.25)--(2.1,-1.45)--(4.9,-1.45)--(4.9,2.25)--(7.8,2.25);
  \draw[deepblue,dashed,very thick] (-.4,-2.25)--(7.8,-2.25);
  \draw[deepblue,very thick,->] (.3,-1.85)--(2.2,-1.85)--(4.0,-1.80)--(5.8,-1.85)--(7.4,-1.85);
  \node[font=\small\bfseries] at (3.5,-2.75) {Group-level continuous partition};
\end{scope}
\end{tikzpicture}
}
\caption{Greedy obstacle-wise side choices make $l^+(s)$ and $l^-(s)$ cross and trigger a fallback stop (a). MIKU's group-level partition keeps a continuous passage for $\mathcal{T}_{\mathrm{ego}}$ (b).}
\label{fig:teaser_en}
\end{figure}

The principal contributions of this work are as follows:
\begin{enumerate}
  \item \textbf{Arrival-time-dependent spatial projection:} Addressing the first gap of time-blind occupancy merging, we formulate an arrival-time-dependent mapping that projects obstacle occupancies into center-coordinate path bounds based on estimated ego arrival times, ensuring the path planner considers only obstacles present when the vehicle traverses each station. This mechanism is developed in Section~\ref{sec:method_projection_en} and isolated by ablation~A1 in Table~\ref{tab:ablation_results}.
  \item \textbf{Group-level continuous lateral partitioning and fast scan:} Addressing the second gap of greedy obstacle-wise deadlocks, we prove the continuous-partition lemma and maximum-gap theorem in Section~\ref{sec:method_band_en}. These theoretical foundations reduce the exponential $2^k$ side assignments for $k$ obstacles to an exact $O(k\log k)$ scan under the fixed-section interval model, verified against exhaustive search across 4,000 deterministic instances and evaluated by ablations~A3--A4 in Table~\ref{tab:ablation_results}.
  \item \textbf{Layered temporal homotopy graph and QP-preserving transfer:} Addressing the third gap of cross-stage constraint loss, we transfer spatial decisions to longitudinal planning via uncertainty-expanded occupancy tubes, safe passage windows, and a layered temporal homotopy graph in Section~\ref{sec:method_temporal_en}. Timing conditions are formulated as linear box constraints that preserve downstream QP structures and computational efficiency, with empirical verification through ablations~A5--A6 in Table~\ref{tab:ablation_results}.
  \item \textbf{Comprehensive empirical and system-level validation:} We systematically evaluate MIKU across 3,500 paired scenarios with rigorous confidence intervals and ablations, validate closed-loop execution in 700 rolling-replanning episodes, and demonstrate software-interface compatibility with the production Apollo 11.0 autonomous driving system in Sections~\ref{sec:experimental_design} and \ref{sec:results}.
\end{enumerate}

Section~\ref{sec:related_work} positions the interface among joint and decomposed planners; Section~\ref{sec:problem} defines the PVD interface and its failure mechanisms; Section~\ref{sec:method} presents MIKU; Sections~\ref{sec:experimental_design}--\ref{sec:results} cover evaluation and results; Sections~\ref{sec:discussion}--\ref{sec:conclusion} discuss and conclude.
